\documentclass[a4paper,12pt]{article}
\usepackage[T2A]{fontenc}
\usepackage[utf8]{inputenc}
\usepackage[russian]{babel}
\usepackage{amsmath,amssymb,mathtools}
\usepackage{helvet}
\renewcommand{\familydefault}{\sfdefault}
\usepackage{geometry}
\geometry{left=22mm,right=22mm,top=23mm,bottom=25mm}
\usepackage{microtype}
\usepackage{array,booktabs,tabularx}
\usepackage{enumitem}
\usepackage{parskip}
\usepackage{xcolor}
\usepackage{tikz}
\usepackage{hyperref}
\definecolor{accent}{HTML}{007C80}
\definecolor{darkaccent}{HTML}{005D65}
\definecolor{deeptext}{HTML}{252B2D}
\definecolor{alert}{HTML}{A52E36}
\hypersetup{colorlinks=true,linkcolor=darkaccent,pdfauthor={Лёвин Артём Александрович},pdftitle={Ревью домашней работы — Ярослав Гаврилов, 29 октября 2026}}
\color{deeptext}
\setlength{\parindent}{0pt}
\setlength{\parskip}{6pt}
\setlist[itemize]{leftmargin=*,itemsep=2pt,topsep=2pt}
\setlength{\emergencystretch}{2em}
\newcommand{\keyidea}{\textcolor{darkaccent}{\textbf{Ключевая идея.}}}
\newcommand{\err}{\textcolor{alert}{\textbf{Ошибка.}}}
\newcommand{\result}{\textcolor{darkaccent}{\textbf{Итог.}}}
\newcommand{\taskheading}[1]{\begin{center}\Large\bfseries Задание № #1\end{center}\vspace{3mm}}
\begin{document}
\thispagestyle{empty}
\begin{center}
\vspace*{29mm}
{\large\color{darkaccent}МАТЕМАТИКА}\par
\vspace{17mm}
{\Huge Ревью домашней работы}\par
\vspace{11mm}
{\LARGE\color{accent}Свойства логарифмов}\par
\vspace{8mm}
{\normalsize Подготовка к ЕГЭ по профильной математике}\par
\vfill
{\large Ярослав Гаврилов}\par
\vspace{14mm}
29 октября 2026 г.\par
\vspace{24mm}
{\small Лёвин Артём Александрович}\par
{\small эксклюзивно для Ярослава Гаврилова}\par
\vspace{14mm}
\begin{tikzpicture}
\draw[accent,line width=0.75pt] (-4,0) .. controls (-2.5,0.62) and (-1.6,-0.54) .. (0,0)
.. controls (1.7,0.55) and (2.6,-0.5) .. (4,0);
\end{tikzpicture}
\end{center}
\newpage
\section*{Сводная таблица}
В таблицу включены только задания, для которых требуется исправление. Нажмите на номер, чтобы перейти к подробному разбору.
\vspace{3mm}
\par\noindent
\renewcommand{\arraystretch}{1.55}
\small
\begin{tabularx}{\dimexpr\linewidth-4pt\relax}{@{}>{\centering\arraybackslash}p{12mm}p{26mm}X>{\raggedright\arraybackslash}p{29mm}>{\raggedright\arraybackslash}p{29mm}@{}}
\toprule
№ & Тема & Суть ошибки & Ваш ответ & Верный ответ\\
\midrule
\hyperref[task:1]{\textbf{1}} & Определение логарифма & Значение логарифма перепутано со степенью числа. & $2^6$ & $6$\\
\hyperref[task:5]{\textbf{5}} & Свойства логарифмов & Не учтён квадрат аргумента при коэффициенте 2. & $\log_2 15$ & $\log_2 75$\\
\hyperref[task:7]{\textbf{7}} & Смена основания и свёртывание & Минус перед логарифмом перенесён на аргумент как знак минус. & $\log_3(-4\cdot5^{3/2})$ & $\log_3\!\dfrac{5\sqrt5}{4}$\\
\bottomrule
\end{tabularx}
\par
\normalsize
\vspace{9mm}
\noindent\textcolor{darkaccent}{\rule{\linewidth}{0.55pt}}
\vspace{5mm}
\textbf{Основной ориентир.} Логарифм показывает показатель степени. При преобразовании суммы и разности логарифмов внимательно проверяйте основания, знаки и коэффициенты перед логарифмами.
\newpage
\thispagestyle{empty}
\vspace*{0.31\textheight}
\begin{center}
{\Huge\bfseries Разбор ошибок}\par
\vspace{8mm}
{\color{darkaccent}\large Три задания требуют исправления}\par
\vspace{11mm}
\begin{tikzpicture}
\draw[accent,line width=0.7pt] (-2.6,0) .. controls (-1.6,0.22) and (-0.8,-0.2) .. (0,0)
.. controls (0.8,0.2) and (1.6,-0.22) .. (2.6,0);
\end{tikzpicture}
\end{center}
\newpage
\phantomsection\label{task:1}
\taskheading{1}
\textbf{Текст задания.} Вычислите $\log_2 64$.

\textbf{Ваше решение.} В работе записано:
\[
\log_2 64=2^6.
\]
\textbf{Ваш ответ.} $2^6$.

\err{} Показатель степени принят за результат возведения в степень.

\textbf{Где именно возникает ошибка.} Правильное наблюдение: $64=2^6$. Первый неверный переход --- равенство $\log_2 64=2^6$. Правая часть этого равенства равна $64$, тогда как логарифм равен $6$.

\textbf{Почему это неверно.} По определению логарифма число $\log_a b$ отвечает на вопрос: «В какую степень нужно возвести основание $a$, чтобы получить $b$?» Здесь основание равно $2$, а число под знаком логарифма равно $64$. Возведение $2$ в шестую степень даёт $64$, поэтому искомый показатель --- $6$. Если записать в ответ $2^6$, получится само число $64$, и определение логарифма будет применено неверно.

\keyidea{} При $a>0$, $a\ne1$ и $b>0$ выполняется
\[
\log_a b=c\quad\Longleftrightarrow\quad a^c=b.
\]
В частности, $\log_a(a^c)=c$.

\textbf{Исправление от последнего правильного шага.} Из представления $64=2^6$ сразу следует
\[
\log_2 64=\log_2(2^6)=6.
\]
\textbf{Полное правильное решение.} Представим число $64$ как степень основания $2$:
\[
64=2\cdot2\cdot2\cdot2\cdot2\cdot2=2^6.
\]
Следовательно, показатель, который требуется найти, равен $6$.

\textbf{Проверка.} $2^6=64$, то есть найденный показатель действительно возвращает число под знаком логарифма.

\result{} $\boxed{6}$.

\vspace{3mm}
\textbf{Задача на закрепление.} Вычислите $\log_3 81$.
\newpage
\phantomsection\label{task:5}
\taskheading{5}
\textbf{Текст задания.} Сверните в один логарифм:
\[\log_2 3+2\log_2 5.\]

\textbf{Ваше решение.} В работе приведено преобразование
\[
\log_2 3+2\log_2 5\quad\longrightarrow\quad\log_2(3\cdot5)=\log_2 15.
\]
\textbf{Ваш ответ.} $\log_2 15$.

\err{} Потерян коэффициент $2$ перед вторым логарифмом.

\textbf{Где именно возникает ошибка.} Исходная сумма записана правильно. Первый неверный шаг --- объединение аргументов $3$ и $5$ без учёта множителя $2$ перед $\log_2 5$.

\textbf{Почему это неверно.} Правило суммы логарифмов применяют после того, как вынесенные перед ними числовые коэффициенты перенесены в показатели степеней аргументов. Множитель $2$ перед $\log_2 5$ означает, что число $5$ следует возвести в квадрат: $5^2=25$. Замена $2\log_2 5$ на $\log_2 5$ меняет значение выражения и приводит к аргументу $15$ вместо $75$.

\keyidea{} При $a>0$, $a\ne1$, $u>0$, $v>0$:
\[
k\log_a u=\log_a(u^k),\qquad
\log_a u+\log_a v=\log_a(uv).
\]

\textbf{Исправление.} Сначала перенесём коэффициент $2$ в показатель степени аргумента второго логарифма:
\[
2\log_2 5=\log_2(5^2)=\log_2 25.
\]
Теперь можно применить правило суммы.

\textbf{Полное правильное решение.}
\begin{align*}
\log_2 3+2\log_2 5
  &=\log_2 3+\log_2(5^2)\\
  &=\log_2 3+\log_2 25\\
  &=\log_2(3\cdot25)\\
  &=\log_2 75.
\end{align*}
\textbf{Проверка.} Раскроем полученный логарифм по свойству произведения:
\[
\log_2 75=\log_2(3\cdot5^2)=\log_2 3+2\log_2 5.
\]
Получили исходное выражение.

\result{} $\boxed{\log_2 75}$.

\vspace{3mm}
\textbf{Задача на закрепление.} Сверните в один логарифм:
$\log_3 2+2\log_3 7$.
\newpage
\phantomsection\label{task:7}
\begingroup
\small
\setlength{\parskip}{3pt}
\setlength{\abovedisplayskip}{5pt}
\setlength{\belowdisplayskip}{5pt}
\setlength{\abovedisplayshortskip}{3pt}
\setlength{\belowdisplayshortskip}{3pt}
\taskheading{7}
\textbf{Текст задания.} Сверните в один логарифм:
\[
\log_{1/3}4+\log_{\sqrt[3]{9}}5.
\]
\textbf{Ваше решение.} Приведение к основанию $3$ выполнено верно:
\[
\log_{1/3}4+\log_{\sqrt[3]{9}}5
=-\log_3 4+\frac32\log_3 5.
\]
Затем в работе отрицательный коэффициент перенесён под знак логарифма как отрицательный аргумент:
\[
\log_3(-4)+\log_3\!\left(5^{3/2}\right)
=\log_3\!\left(-4\cdot5^{3/2}\right).
\]
\textbf{Ваш ответ.} $\log_3(-4\cdot5^{3/2})$.

\err{} Неправильный перенос минуса под знак логарифма.

\textbf{Где именно возникает ошибка.} Последний правильный шаг --- $-\log_3 4+\tfrac32\log_3 5$. Первый неверный шаг --- замена $-\log_3 4$ на $\log_3(-4)$.

\textbf{Почему это неверно.} При умножении логарифма на $-1$ изменяется показатель степени его положительного аргумента: $-\log_3 4=\log_3(4^{-1})=\log_3\tfrac14$. Число $-4$ не подходит в качестве аргумента действительного логарифма: аргумент должен быть строго положительным. Поэтому запись $\log_3(-4)$ вообще не определена в действительных числах.

\keyidea{} Для положительного $u$ верно
\[
-\log_a u=\log_a\frac1u,
\qquad\frac32\log_a u=\log_a u^{3/2}.
\]
При свёртывании логарифмов аргументы должны оставаться положительными.

\textbf{Исправление от последнего правильного шага.}
\begin{align*}
-\log_3 4+\frac32\log_3 5
 &=\log_3\frac14+\log_3\left(5^{3/2}\right)\\
 &=\log_3\left(\frac{5^{3/2}}4\right).
\end{align*}
\textbf{Полное правильное решение.} Приведём оба основания к степеням числа $3$:
\[
\frac13=3^{-1},\qquad \sqrt[3]{9}=3^{2/3}.
\]
Используем правило смены степени основания $\log_{3^r}u=\frac1r\log_3 u$ при $r\ne0$. Тогда
\begin{align*}
\log_{1/3}4+\log_{\sqrt[3]{9}}5
 &=\frac1{-1}\log_3 4+\frac1{2/3}\log_3 5\\
 &=-\log_3 4+\frac32\log_3 5\\
 &=\log_3\frac14+\log_3\left(5\sqrt5\right)\\
 &=\log_3\frac{5\sqrt5}{4}.
\end{align*}
\textbf{Проверка.} Аргумент $\dfrac{5\sqrt5}{4}$ положителен, основание $3>0$ и $3\ne1$. Раскрыв логарифм частного, вновь получаем $-\log_3 4+\tfrac32\log_3 5$.

\result{} $\boxed{\displaystyle\log_3\frac{5\sqrt5}{4}}$.

\vspace{3mm}
\textbf{Задача на закрепление.} Сверните в один логарифм:
$\log_{1/4}9+\log_{\sqrt[3]{8}}5$.
\endgroup
\end{document}
